\documentclass[a4paper,11pt]{article}
\usepackage[utf8]{inputenc}
\usepackage{amsmath, amssymb, amsfonts}
\usepackage{bm}
\usepackage{graphicx}
\usepackage{geometry}
\usepackage{enumitem}
\usepackage{subcaption}
\usepackage{float}
\geometry{margin=2.5cm}

% \title{Fisher-Kolmogorov equation for neurodegenerative diseases}
% \author{Based on handwritten notes by Ema and reconstructed in LaTeX}
% \date{}

\title{\huge \textbf{Fisher-Kolmogorov Equation} \\[0.2em] 
       \huge \textbf{A Mathematical Model for Neurodegenerative Diseases} \\[0.8em]
       \Large Numerical Analysis for PDE Project}
\author{Emanuele Severino, Nicola Noventa, Matteo Arrigo}
\date{\today}

\begin{document}
\maketitle

% Generate the Table of Contents
\tableofcontents
\newpage  % Start content on a new page

\section{Introduction}
Create a model that explains the growth and spreading of misfolded proteins. This type of modeling is particularly relevant in the context of neurodegenerative diseases, where the accumulation and propagation of misfolded proteins can trigger pathological cascades.

\section{Mathematical Modelling}
\subsection{Kinetics of Prion-like Spreading}
Let:
\begin{itemize}
    \item $\rho$ be the concentration of healthy proteins.
    \item $\tilde{\rho}$ be the concentration of misfolded proteins.
\end{itemize}

We consider the protein conversion reaction:
\begin{equation}
    \rho + \tilde{\rho} \xrightarrow{k_{12}} \tilde{\rho} + \tilde{\rho}
\end{equation}
This autocatalytic reaction models the idea that a misfolded protein can convert a healthy one upon interaction.

\paragraph{Parameter:}
\begin{itemize}
    \item $k_{12}$: rate of conversion from healthy to misfolded proteins, proportional to both $\rho$ and $\tilde{\rho}$.
\end{itemize}

\subsection{Governing Equations}
The spatio-temporal evolution of the concentrations is modeled as:
\begin{align}
    \frac{d \rho}{dt} &= \nabla \cdot (D_{\rho} \nabla \rho) + k_0 - k_1 \rho - k_{12} \rho \tilde{\rho} \\
    \frac{d \tilde{\rho}}{dt} &= \nabla \cdot (D_{\tilde{\rho}} \nabla \tilde{\rho}) + k_{12} \rho \tilde{\rho} - \tilde{k}_1 \tilde{\rho}
\end{align}
These equations include production, clearance, diffusion, and nonlinear reaction terms.

\paragraph{Where:}
\begin{itemize}
    \item $k_0$: production rate of healthy protein $\rho$ (e.g., from synthesis).
    \item $k_1$: clearance rate of $\rho$ (e.g., degradation or dilution).
    \item $\tilde{k}_1$: clearance rate of $\tilde{\rho}$.
    \item $k_{12}$: nonlinear conversion rate.
    \item $D_{\rho}, D_{\tilde{\rho}}$: diffusion coefficients that characterize the spatial spread.
\end{itemize}

\subsection{Stationary States}
\paragraph{Healthy state ($\tilde{\rho} = 0$):}
No misfolded proteins are present, and the system reaches equilibrium for $\rho$:
\begin{equation}
    \frac{d \rho}{dt} = 0 \Rightarrow k_0 - k_1 \rho = 0 \Rightarrow \rho = \frac{k_0}{k_1}
\end{equation}
This is the baseline equilibrium in the absence of pathological proteins.

\paragraph{Initial state ($\tilde{\rho} \ll \rho$):}
When misfolded proteins are rare, we can assume a quasi-steady state for $\rho$:
\begin{equation}
    \frac{d \rho}{dt} \approx 0, \quad \nabla \cdot (D_{\rho} \nabla \rho) \approx 0
\end{equation}
Leading to:
\begin{equation}
    k_0 - k_1 \rho - k_{12} \rho \tilde{\rho} = 0 \Rightarrow \rho = \frac{k_0}{k_1 + k_{12} \tilde{\rho}}
\end{equation}
This gives a rational expression for $\rho(\tilde{\rho})$ which reflects the feedback inhibition induced by $\tilde{\rho}$.

\paragraph{Linear Approximation:}
If $\tilde{\rho} \ll 1$, then using a Taylor expansion:
\begin{equation}
    \frac{1}{k_1 + k_{12} \tilde{\rho}} \approx \frac{1}{k_1} \left(1 - \frac{k_{12}}{k_1} \tilde{\rho} \right)
\end{equation}
Therefore,
\begin{equation}
    \rho \approx \frac{k_0}{k_1} \left(1 - \frac{k_{12}}{k_1} \tilde{\rho} \right)
\end{equation}
This approximation will be useful for simplifying the equations governing $\tilde{\rho}$.

\subsection{Reduced Equation and Normalization}
Substitute the linear approximation of $\rho$ into the equation for $\tilde{\rho}$:
\begin{align}
    \frac{d \tilde{\rho}}{dt} &= \nabla \cdot (D_{\tilde{\rho}} \nabla \tilde{\rho}) + \left( \frac{k_{12} k_0}{k_1} - \tilde{k}_1 \right) \tilde{\rho} - \frac{k_{12}^2 k_0}{k_1^2} \tilde{\rho}^2
\end{align}
This is a reaction-diffusion equation with logistic-type growth and saturation.

\paragraph{Normalization:}
Define the maximum misfolded concentration $\tilde{\rho}_{\max}$ and the relative concentration:
\begin{equation}
    c = \frac{\tilde{\rho}}{\tilde{\rho}_{\max}}, \quad 0 \leq c \leq 1
\end{equation}
Then the equation becomes:
\begin{equation}
    \frac{\partial c}{\partial t} = \nabla \cdot (D \nabla c) + \alpha c (1 - c), \quad \text{with } \alpha = \frac{k_{12} k_0}{k_1} - \tilde{k}_1
\end{equation}
This is the classical Fisher-KPP equation, modeling the propagation of concentration fronts.

\subsection{Steady Solutions}
Setting \( \partial c/\partial t = 0 \), we obtain the stationary form:
\begin{equation}
    \nabla \cdot (D \nabla c) + \alpha c(1 - c) = 0
\end{equation}
This admits equilibrium solutions $c = 0$ (unstable) and $c = 1$ (stable), which represent the absence or saturation of misfolded proteins in the tissue.

\subsection{Diffusion Term Details}
\begin{itemize}
    \item \textbf{Isotropic diffusion:} the diffusion tensor $D$ is scalar, uniform in all directions:
    \begin{equation}
        \nabla \cdot (D \nabla c)
    \end{equation}
    \item \textbf{Anisotropic diffusion:} maximal along preferred direction $\bm{m}$:
    \begin{equation}
        D = D_{\text{const}} (I + d_{\text{an}} \bm{m} \otimes \bm{m})
    \end{equation}
    and the diffusion term becomes:
    \begin{equation}
        \nabla \cdot (D \nabla c) = \nabla \cdot \left[ D_{\text{const}} (I + d_{\text{an}} \bm{m} \otimes \bm{m}) \nabla c \right]
    \end{equation}
\end{itemize}
This form accounts for directional spreading influenced by structural anisotropies (e.g., in tissue).

\section{Mathematical Model: Generalized Fisher–Kolmogorov Equation}

The proposed model is a nonlinear anisotropic reaction--diffusion equation that describes the evolution of the concentration \( c = c(\mathbf{x}, t) \) of misfolded proteins over space and time. It reads as follows:

\begin{equation}
\begin{cases}
\displaystyle \frac{\partial c}{\partial t} - \nabla \cdot (D \nabla c) - \alpha c (1 - c) = 0 & \text{in } \Omega, \\
D \nabla c \cdot \mathbf{n} = 0 & \text{on } \partial \Omega, \\
c(t=0) = c_0 & \text{in } \Omega.
\end{cases}
\label{eq:fisher-kolmogorov}
\end{equation}

\subsection{Weak Formulation}

To derive the weak formulation of the Fisher--Kolmogorov equation, we begin by introducing the appropriate functional setting. The natural Hilbert space for both the solution and the test functions is defined as:
\[
V := H^1(\Omega) = \left\{ v \in L^2(\Omega) \, : \, \nabla v \in \left[L^2(\Omega)\right]^d \right\},
\]
where \( L^2(\Omega) \) denotes the space of square-integrable functions over the domain \( \Omega \), and \( d \) is the dimension of the spatial domain.

\medskip

\noindent
\textbf{Step 1: Multiplication by a test function and integration over the domain.}

\smallskip

Let \( v \in H^1(\Omega) \) be a test function. Multiplying the strong form of the equation by \( v \) and integrating over \( \Omega \), we obtain:
\[
\int_\Omega \frac{\partial c}{\partial t} v \, dx 
- \int_\Omega \nabla \cdot (D \nabla c) \, v \, dx 
- \int_\Omega \alpha c(1 - c) v \, dx = 0
\qquad \forall v \in H^1(\Omega).
\]

\medskip

\noindent
\textbf{Step 2: Integration by parts of the diffusion term.}

\smallskip

The second term, involving the divergence of the flux, is integrated by parts:
\[
- \int_\Omega \nabla \cdot (D \nabla c) \, v \, dx 
= - \int_{\partial \Omega} (D \nabla c \cdot \mathbf{n}) v \, d\sigma 
+ \int_\Omega D \nabla c \cdot \nabla v \, dx.
\]

\noindent
Using the homogeneous Neumann boundary condition
\[
D \nabla c \cdot \mathbf{n} = 0 \quad \text{on } \partial \Omega,
\]
the boundary integral vanishes, leaving only the volume term.

\medskip

\noindent
\textbf{Final weak formulation.}

\smallskip

We can now state the weak formulation as follows:

Find \( c \in H^1(\Omega) \) such that
\begin{equation}
\int_\Omega \frac{\partial c}{\partial t} \, v \, dx 
+ \int_\Omega D \nabla c \cdot \nabla v \, dx 
- \int_\Omega \alpha c(1 - c) v \, dx = 0, 
\qquad \forall v \in H^1(\Omega)
\label{eq:weak-continuous}
\end{equation}

\subsection{Time Discretization: Backward Euler Scheme}

To numerically approximate the solution of the time-dependent weak formulation, we begin by discretizing the temporal domain and constructing a sequence of variational problems at discrete time levels.

\subsubsection*{Temporal partitioning}

Let \( T > 0 \) be the final simulation time. We partition the time interval \( [0, T] \) uniformly into \( N \) subintervals of equal length \( \Delta t = \frac{T}{N} \), and define the discrete time instants as
\[
t^n := n \Delta t, \qquad \text{for } n = 0, 1, \dots, N
\]

At each time level \( t^n \), we aim to compute an approximation \( c^n \approx c(t^n) \) of the exact solution \( c(x, t^n) \).

\subsubsection*{Approximation of the time derivative}

The time derivative \( \frac{\partial c}{\partial t} \) is approximated using the backward (implicit) Euler scheme, evaluated at time \( t^{n+1} \):
\[
\left. \frac{\partial c}{\partial t} \right|_{t^{n+1}} \approx \frac{c^{n+1} - c^n}{\Delta t} 
\]

This transforms the time-dependent problem into a sequence of quasi-static (elliptic-type) problems. At each time step, we compute the unknown \( c^{n+1} \) given the known solution at the previous time step \( c^n \).

\subsubsection*{Fully discrete variational problem}

Starting from the weak formulation of the continuous problem \eqref{eq:weak-continuous} we replace the time derivative by using the backward Euler approximation.

The resulting time-discrete variational problem reads:


\textbf{Given} \( c^n \in H^1(\Omega) \), \textbf{find} \( c^{n+1} \in H^1(\Omega) \) such that:
\begin{equation}
\int_\Omega \frac{c^{n+1} - c^n}{\Delta t} \, v \, dx
+ \int_\Omega D \nabla c^{n+1} \cdot \nabla v \, dx
- \int_\Omega \alpha c^{n+1}(1 - c^{n+1}) v \, dx = 0,
\qquad \forall v \in H^1(\Omega)
\label{eq:BE_applied}
\end{equation}

\medskip

This problem is nonlinear in \( c^{n+1} \) due to the reaction term \( \alpha c^{n+1}(1 - c^{n+1}) \). It must be solved at each time step, starting from the initial condition \( c^0 = c_0 \in H^1(\Omega) \).

\subsection{Linearization via Fréchet Derivative}

% Given the nonlinear variational problem derived from the backward Euler time discretization:

% \begin{quote}
% \textbf{Given} \( c^n \in H^1(\Omega) \), \textbf{find} \( c^{n+1} \in H^1(\Omega) \) such that
% \[
% \int_\Omega \frac{c^{n+1} - c^n}{\Delta t} v \, dx + \int_\Omega D \nabla c^{n+1} \cdot \nabla v \, dx - \int_\Omega \alpha c^{n+1}(1 - c^{n+1}) v \, dx = 0 \qquad \forall v \in H^1(\Omega),
% \]
% \end{quote}

Given \eqref{eq:BE_applied} we define the associated \textbf{residual functional} \( R \colon H^1(\Omega) \to H^1(\Omega)' \) by:
\[
R(c, v) := \int_\Omega \frac{c - c^n}{\Delta t} v \, dx + \int_\Omega D \nabla c \cdot \nabla v \, dx - \int_\Omega \alpha c(1 - c) v \, dx.
\]

Solving the variational problem is equivalent to finding \( c^{n+1} \in H^1(\Omega) \) such that:
\[
R(c^{n+1}, v) = 0 \qquad \forall v \in H^1(\Omega).
\]

\subsubsection*{Fréchet derivative: definition}

To solve the nonlinear problem, we linearize the residual \( R \) via its \textbf{Fréchet derivative}. Let \( \delta \in H^1(\Omega) \) be a perturbation direction. The Fréchet derivative of \( R \) at a point \( c \in H^1(\Omega) \) in the direction \( \delta \) is defined as:
\[
R'(c)(\delta, v) := \left. \frac{d}{d\varepsilon} R(c + \varepsilon \delta, v) \right|_{\varepsilon = 0}.
\]

This way it can be seen that holds
\begin{equation}
R(c + \delta, v) = R(c, v) + R'(c)(\delta, v)
\qquad as \ \delta \rightarrow 0
\label{frechet_approx}
\end{equation}

\subsubsection*{Computation of the derivative}

We expand the residual functional evaluated at a perturbed function \( c + \delta \):
\[
R(c + \delta, v) = \int_\Omega \frac{c + \delta - c^n}{\Delta t} v \, dx + \int_\Omega D \nabla(c + \delta) \cdot \nabla v \, dx - \int_\Omega \alpha (c + \delta)(1 - c - \delta) v \, dx.
\]

Focusing on the nonlinear reaction term:
\[
(c + \delta)(1 - c - \delta) = c(1 - c) - \delta(1 - 2c) - \delta^2.
\]

Substituting into the residual:
\begin{align*}
R(c + \delta, v) &= \int_\Omega \left( \frac{c - c^n}{\Delta t} + \frac{\delta}{\Delta t} \right) v \, dx 
+ \int_\Omega D (\nabla c + \nabla \delta) \cdot \nabla v \, dx \\
&\quad - \int_\Omega \alpha \left[ c(1 - c) - \delta(1 - 2c) - \delta^2 \right] v \, dx.
\end{align*}

Grouping terms:
\begin{align*}
R(c + \delta, v) &= R(c, v) + \int_\Omega \frac{\delta}{\Delta t} v \, dx + \int_\Omega D \nabla \delta \cdot \nabla v \, dx - \int_\Omega \alpha (1 - 2c) \delta v \, dx + \mathcal{O}(\|\delta\|^2).
\end{align*}

Taking \( \delta \in H^1(\Omega) \) so that \( \delta \rightarrow 0\), the remainder \( \mathcal{O}(\|\delta\|^2) \) can be neglected in the linear approximation. Therefore, comparing with \eqref{frechet_approx}, the Fréchet derivative is:

\[
R'(c)(\delta, v) = \int_\Omega \frac{\delta}{\Delta t} v \, dx + \int_\Omega D \nabla \delta \cdot \nabla v \, dx - \int_\Omega \alpha (1 - 2c) \delta v \, dx.
\]

\subsection{Application of Newton's method}

To solve the nonlinear problem \( R(c^{n+1}, v) = 0 \), we apply the Newton method. Starting from an initial guess \( c^{(0)} \), we iterate:

\begin{quote}
\textbf{At each iteration} \( k \), find \( \delta^{(k)} \in H^1(\Omega) \) such that
\[
R'(c^{(k)})(\delta^{(k)}, v) = -R(c^{(k)}, v) \qquad \forall v \in H^1(\Omega),
\]
and update the approximation:
\[
c^{(k+1)} = c^{(k)} + \delta^{(k)}.
\]
\end{quote}

The iteration is repeated until convergence, e.g., when \( \|\delta^{(k)}\| \) or \( \|R(c^{(k)}, \cdot)\|_{H^1(\Omega)'} \) is below a fixed tolerance.

\subsubsection{Finite Element Discretization and Assembly of the Newton System}

We now discretize the problem in space using the Galerkin finite element method. Let \( V_h \subset H^1(\Omega) \) be a finite-dimensional subspace spanned by a basis \( \{\phi_j\}_{j=1}^N \). We approximate \( c^{(k)} \) and \( \delta^{(k)} \) as:

\[
c^{(k)}(x) = \sum_{j=1}^N c_j^{(k)} \phi_j(x), \qquad \delta^{(k)}(x) = \sum_{j=1}^N \delta_j^{(k)} \phi_j(x).
\]

Choosing the test functions \( v = \phi_i \), for \( i = 1, \dots, N \), we project the weak formulation onto the space \( V_h \), resulting in a linear system of the form:

\[
\mathbf{A}^{(k)} \delta^{(k)} = -\mathbf{R}^{(k)},
\]

where:

\begin{itemize}
  \item \( \delta^{(k)} = (\delta_1^{(k)}, \dots, \delta_N^{(k)})^T \) is the vector of unknown Newton corrections,
  \item \( \mathbf{A}^{(k)} \in \mathbb{R}^{N \times N} \) is the Jacobian matrix (the discretization of the Fréchet derivative),
  \item \( \mathbf{R}^{(k)} \in \mathbb{R}^N \) is the residual vector evaluated at \( c^{(k)} \).
\end{itemize}

The entries of the matrix \( \mathbf{A}^{(k)} \) are given by:

\[
A_{ij}^{(k)} = \int_\Omega \frac{1}{\Delta t} \phi_j \phi_i \, dx 
+ \int_\Omega D \nabla \phi_j \cdot \nabla \phi_i \, dx 
- \int_\Omega \alpha (1 - 2c^{(k)}) \phi_j \phi_i \, dx,
\]

and the residual vector by:

\[
R_i^{(k)} = \int_\Omega \frac{c^{(k)} - c^n}{\Delta t} \phi_i \, dx 
+ \int_\Omega D \nabla c^{(k)} \cdot \nabla \phi_i \, dx 
- \int_\Omega \alpha c^{(k)} (1 - c^{(k)}) \phi_i \, dx.
\]

At each Newton iteration \( k \), the linear system \( \mathbf{A}^{(k)} \delta^{(k)} = -\mathbf{R}^{(k)} \) is solved, and the solution is updated:

\[
c^{(k+1)} = c^{(k)} + \delta^{(k)}.
\]



\section{Computational Methods}

\subsection{Computational Problem Formulation}
To sum up, the problem is formalized as the Fisher–Kolmogorov Equation, a specific nonlinear partial differential equation (PDE) over a domain in \(\mathbb{R}^3\) representing the brain. The boundary conditions are homogeneous Neumann conditions on the entire surface of the 3D domain, and the initial condition is provided.

The solver starts from the initial solution and iterates over time until reaching a specified endpoint, using a discretization determined by a time step. At each time step, our software solves the linear system corresponding to the linearized weak formulation of the PDE, discretized in space using the finite element method (FEM). Specifically, the solver assembles the system matrix and the right-hand side vector and then solves the system using a linear solver.

We construct the initial condition as a function with a nonzero constant value (initial seed) only within a defined region (initial seed region). To simplify the modeling, we support an initial seed region that is cubic in shape (and intersected with the domain), with a fixed edge length.

\subsection{Simulation Setup}

The inputs and parameters of the problem are:
\begin{enumerate}
    \item \textbf{Numerical Parameters}
    \begin{enumerate}[label=-]
        \item \textbf{T}: Final time of the simulation.
        \item \textbf{\( \Delta t \)}: Time step for time discretization.
        \item \textbf{Mesh}: A 3D discretized representation of the domain in \(\mathbb{R}^3\).
        \item \textbf{Degree}: The polynomial degree of the finite elements used to construct the FEM space over the mesh.
    \end{enumerate}
    \item \textbf{Biological Parameters}
    \begin{enumerate}[label=-]
        \item \textbf{d\_ext}: Isotropic extracellular diffusion coefficient.
        \item \textbf{d\_axn}: Anisotropic axonal transport along axonal fiber directions.
        \item \textbf{Diffusion Type}: Defines the predefined axonal fiber direction to be considered.
        \begin{itemize}
            \item Isotropic
            \item Circumferential
            \item Radial
        \end{itemize}
        \item \textbf{$\alpha$}: Growth rate of misfolded protein concentration.
        \item \textbf{Initial Seed}: Value of the initial nonzero concentration.
        \item \textbf{Seed Point}: Center of the initial seed region.
    \end{enumerate}
\end{enumerate}
All parameters are assumed to remain constant over time.

\textbf{Outputs}: The solver outputs \texttt{.vtk} and \texttt{.pvtk} files representing the mesh with the computed solution at each time step. These results can be visualized using ParaView.

The primary objects in the implementation include:
\begin{itemize}[label=-, itemsep=0pt, parsep=5pt]
    \item \textbf{Residual Vector}: Corresponds to \( R \), representing the right-hand side (RHS) of the linearized system.
    \item \textbf{Jacobian Matrix}: Corresponds to \( A \), representing the left-hand side (LHS) matrix of the linearized system.
    \item \textbf{Delta}: Corresponds to \( \delta \), representing the unknown increment in the linearized system.
    \item \textbf{Solution}: Corresponds to \( c \), representing the global solution of the problem.
\end{itemize}

\subsection{Implementation of the 1D Solver}

To validate our approach, we first implemented the solver for a corresponding 1D problem, as the authors of the reference paper did. In this case, the problem parameters are simplified:
\begin{itemize}[itemsep=0pt, parsep=5pt]
    \item The domain is 1D: we consider the interval \([-1,1] \subset \mathbb{R}\) as the fixed domain.
    \item The diffusion term \( D \) reduces to a scalar parameter \( d \).
    \item The initial seed region consists of a single point at the center of the domain (\( x = 0 \)).
\end{itemize}

The structure of the code follows the template used in our laboratory sessions. The mesh is created in the given interval with \( N = 200 \) mesh elements and \( \text{degree} = 1 \), resulting in \( 201 \) total degrees of freedom (DoFs). After initializing all necessary objects, the method \texttt{solve()} executes the system resolution for each time step and writes the output for that instant.

The method \texttt{solve()} calls \texttt{solve\_newton()}, which performs Newton iterations. For each iteration, the system is linearized around the current solution via \texttt{assemble\_system()}, and the resulting linear system is solved using \texttt{solve\_linear\_system()}. This process repeats until convergence is achieved, which we define as when the norm of the residual vector falls below a fixed threshold (\( 10^{-6} \) in our code).

\textbf{Assemble\_system()} follows the structure of element-wise loops used in the laboratory exercises to compute all contributions to the three integral terms in the LHS matrix (for each pair of FEM basis functions) and the three integral terms in the RHS vector (for each basis function). These terms represent the linearized versions of the original system, which we obtain using the values of the current and previous solutions to compute the relevant quantities.

Since homogeneous Neumann boundary conditions introduce no additional terms in the mathematical formulation, and no Dirichlet conditions are imposed, the code does not explicitly handle boundary conditions.

\textbf{Solve\_linear\_system()} assumes that the LHS matrix and RHS vector are correctly assembled. Since the problem is symmetric, we use the Conjugate Gradient (CG) method to numerically solve the resulting linear system. Moreover, an SSOR preconditioner is applied to accelerate convergence.

\subsection{Implementation of the 3D Solver}

The primary differences in the 3D solver involve parallelization, data representation, and the linear solver.

\subsubsection{Parallelization Strategy}

The computational complexity of the 3D solver is significantly higher due to the large number of nodes in the 3D mesh, the nonlinearity of the problem, and the large number of time steps required. Therefore, parallelization is necessary to accelerate simulations and make them computationally feasible.

Our parallelization strategy follows the approach covered in our course. Many objects are converted into their parallel equivalents, particularly those related to linear algebra, which are now wrapped using Trilinos. Trilinos facilitates parallel computations using MPI. The mesh is partitioned among MPI processes, with each process handling its owned and relevant DoFs while ensuring correct communication between DoFs at partition boundaries. The solution and delta vectors are managed similarly, abstracting much of the parallel complexity.

\subsubsection{Data Representation and Mesh Handling}
In 3D, data representation is more complex, as we now deal with functions over a three-dimensional domain with a nontrivial geometry. The mesh is read from an external file rather than being generated within the code. We processed the mesh provided alongside the reference paper and converted it into a valid \texttt{.msh} file.

Due to the absence of real-world data and for simplicity, we did not distinguish between grey and white matter regions in the brain.

\textbf{FunctionD} computes the \( 3 \times 3 \) diffusion matrix \( D \) for each node in the mesh using predefined models. The normal vector is obtained through \texttt{compute\_normal\_vector(const Point<dim> \&p)}, following three possible models of fiber orientation. Due to the lack of real axonal fiber data, we could not implement a biologically accurate fiber orientation model.

\textbf{FunctionAlpha} remains a constant function throughout the domain.

\textbf{FunctionU0} computes the initial seed region as the intersection of a cubic area centered at the given seed point and the domain mesh.

\subsubsection{Changes in the Solver}

Since the system matrix in 3D is not guaranteed to be symmetric, we switched from the CG method to the GMRES method, which is more robust for non-symmetric systems. Consequently, we also replaced the SSOR preconditioner with an ILU preconditioner.



\section{Stability, Convergence}

\section{Benchmark}

\subsection{Introduction}
This section presents the benchmarking results of the Fisher-Kolmogorov solver, with a focus on parallel performance. The objective is to evaluate the scalability of the implementation and the impact of parallelization on execution time.


The analysis was conducted on two different platforms:
\begin{itemize}
    \item A local workstation, to evaluate the solver's behaviour in a standard desktop environment.
    \item The MareNostrum 5 supercomputer, to assess the scalability on a high-performance computing system.
\end{itemize}

\subsection{Benchmarking Setup}

The benchmarking tests were conducted on two platforms with different hardware configurations, summarized in Table~\ref{tab:hardware}. On the local workstation, the solver was executed within a WSL2 virtualized environment, while on MareNostrum 5 the tests were run on dedicated compute nodes.

\begin{table}[h]
    \centering
    \begin{tabular}{|c|c|c|c|c|}
        \hline
        Platform & CPU Model & Cores & RAM per Node & Network \\
        \hline
        Local workstation & Intel Core i7-1165G7 & 4 & 8 GB & Shared memory \\
        MareNostrum 5 & 2 × Xeon 8480+ (Sapphire Rapids) & 112 & 256 GB & NDR200 (100 Gb/s) \\
        \hline
    \end{tabular}
    \caption{Hardware configuration of the platforms used for benchmarking.}
    \label{tab:hardware}
\end{table}


The problem size was kept constant for all tests, with a mesh of 634,472 elements and 111,951 degrees of freedom. The simulation covered a total physical time of \( T = 10 \) with a time step size of \( \Delta t = 0.1 \), resulting in 100 time steps.

The number of MPI processes used during the benchmarking was:
\begin{itemize}
    \item Local workstation: 1, 2, and 4 tasks.
    \item MareNostrum 5: from 1 up to 256 tasks.
\end{itemize}


\subsection{Measurement Methodology}

The performance analysis was based on three metrics, each measured separately:

\begin{itemize}
    \item \textit{Setup time:} time required for mesh loading and partitioning, finite element space initialization, degree of freedom distribution, and matrix and vector allocation. This phase corresponds to the execution of the \texttt{setup()} function.
    
    \item \textit{Solve time:} time needed to solve the nonlinear system over all time steps. It includes the application of the initial condition, assembly of system matrices and residuals at each time step, iterative solution of the linear systems, and output writing. It corresponds to the execution of the \texttt{solve()} function.
    
    \item \textit{Total time:} overall execution time, including setup and solve phases, as well as MPI initialization and finalization.
\end{itemize}

Timing data were collected using \texttt{MPI\_Wtime()} and recorded in CSV files for post-processing and analysis.

\subsection{Results and Discussion}

\subsubsection{Local Execution}

The following figures show the setup time, solve time, and total execution time measured on the local workstation as a function of the number of MPI tasks. As expected, the solve time decreases when increasing the number of tasks, while the setup time remains approximately constant. The total time follows the same trend as the solve time, as it represents the dominant contribution.

\begin{figure}[h]
    \centering
    \begin{subfigure}{0.32\textwidth}
        \includegraphics[width=\linewidth]{image/setuplocal.png}
        \caption{Setup time}
    \end{subfigure}
    \hfill
    \begin{subfigure}{0.32\textwidth}
        \includegraphics[width=\linewidth]{image/solvelocal.png}
        \caption{Solve time}
    \end{subfigure}
    \hfill
    \begin{subfigure}{0.32\textwidth}
        \includegraphics[width=\linewidth]{image/totallocal.png}
        \caption{Total time}
    \end{subfigure}
    \caption{Execution times on the local workstation.}
    \label{fig:performance_local}
\end{figure}


The following figures report the setup, solve, and total execution times measured on MareNostrum 5. The results show a good scalability up to 64 tasks, beyond which the communication overhead starts to degrade the parallel efficiency. In particular, the solve time does not decrease significantly when increasing the number of tasks beyond this point, and in some cases it increases due to the cost of inter-node communication.

\begin{figure}[H]
    \centering
    \makebox[\linewidth]{ % Forza la tabella a essere centrata
        \begin{tabular}{c c}
            \includegraphics[width=0.55\textwidth]{image/output.png} &
            \includegraphics[width=0.55\textwidth]{image/output3.png} \\
            \includegraphics[width=0.55\textwidth]{image/output1.png} &
            \includegraphics[width=0.55\textwidth]{image/output4.png} \\
            \includegraphics[width=0.55\textwidth]{image/output2.png} &
            \includegraphics[width=0.55\textwidth]{image/output5.png} \\
        \end{tabular}
    }
    \caption{Comparison of execution times with linear (left) and log-log (right) scales for setup time (top), solve time (middle), and total time (bottom).}
    \label{fig:performance}
\end{figure}

\subsection{Scalability Metrics}

To further evaluate the parallel performance of the solver, two classical scalability metrics were computed: the \textit{speedup} and the \textit{parallel efficiency}.

The \textbf{speedup} \( SU_p \) is defined as the ratio between the execution time on a single processor and the execution time on \( p \) processors:
\begin{equation}
SU_p(n) = \frac{T_1(n)}{T_p(n)}
\end{equation}
where \( T_1(n) \) is the total execution time with one process, and \( T_p(n) \) is the total execution time with \( p \) processes.

The \textbf{parallel efficiency} \( E_p \) is defined as the ratio between the speedup and the number of processors used:
\begin{equation}
E_p(n) = \frac{SU_p(n)}{p} = \frac{T_1(n)}{p \times T_p(n)}
\end{equation}

These metrics provide a quantitative measure of how effectively the solver utilizes the available computational resources.

\subsubsection{Results}

The following table summarizes the speedup and efficiency computed for both the local workstation and MareNostrum 5:


The results confirm that the solver achieves good scalability on MareNostrum 5 up to 64 tasks, with an efficiency close to ideal. Beyond this point, efficiency decreases due to increased communication overhead. On the local workstation, the scalability is limited by the low number of physical cores and the virtualization overhead introduced by WSL2.


\end{document}